% ==============================================================================
%  CHAPTER 4 — RESULTS
% ==============================================================================
\section{Results}
\label{sec:results}

This chapter presents the empirical evaluation of the proposed context-aware
privacy framework. Section~4.1 describes the evaluation dataset and ground-truth
annotation protocol. Section~4.2 reports classification performance for the
geometry-only (Phase~2a) configuration. Section~4.3 extends this to the full
system that additionally incorporates facial landmark cues (Phase~2b), evaluated
first on still frames. Section~4.4 reports a second, larger evaluation of Phase~2b
using short video clips with sound. Section~4.5 reports an offline second pass for
recorded video that preserves confirmed partners and their group for their whole
time on screen. Section~4.6 compares all configurations against the
blanket-blurring baseline. Section~4.7 reports computational performance on the
evaluation hardware and provides a projected estimate for the Raspberry~Pi~5
deployment target.

\subsection{Dataset and Ground-Truth Annotation}
\label{sec:dataset}

The evaluation dataset comprises 21 egocentric video clips that I collected from
publicly available online sources, including YouTube; I did not record any of the
footage myself. All clips are filmed from a first-person viewpoint, with the camera worn
or held at approximately eye level. The clips are divided into two categories:
14 \emph{interaction} clips, in which the camera wearer conducts a face-to-face
conversation with at least one person, and 7 \emph{bystander} clips, in which no
direct interaction takes place and all visible faces belong to people who are
present but not conversing with the wearer.

The two categories differ in setting. The bystander clips show the wearer walking
through busy outdoor public spaces, such as street markets and city streets, by day
and at night, and range from approximately 6 to 78 seconds in length. The
interaction clips are five-minute segments filmed indoors, mostly showing the wearer
playing card or board games with others around a table, with the remainder showing
the wearer shopping in small stores. Source resolutions range from
$1280 \times 720$ to $3840 \times 2160$; all frames were resized to
$1280 \times 720$ for processing. For the Phase~2a evaluation, I manually
labelled 66 face tracks across the 21 clips.

I established the ground-truth labels by watching each clip in full. A track was assigned the label \emph{interaction partner} (1) if the
corresponding person engaged in face-to-face conversation with the camera wearer at
any point during the clip, and \emph{bystander} (0) otherwise. This per-track
binary protocol is consistent with the interaction latching behaviour of the classifier,
which permanently commits a track to the \emph{Interactor} state once the
interaction criteria are met.

Across all 66 tracks, 26 were labelled as interaction partners and 40 as bystanders
(Table~\ref{tab:dataset}).

\begin{table}[htbp]
\begin{singlespace}
\centering
\caption{Evaluation dataset composition and ground-truth track distribution.}
\label{tab:dataset}
\renewcommand{\arraystretch}{1.3}
\begin{tabular}{@{}lc@{}}
\toprule
\textbf{Video clips} & \textbf{Count} \\
\midrule
Interaction clips (5-minute indoor segments) & 14 \\
Bystander clips (6-78~s outdoor walking)    & 7  \\
\textbf{Total clips}                         & \textbf{21} \\
\midrule
\textbf{Labelled face tracks} & \textbf{Count} \\
\midrule
Interaction partner tracks & 26 \\
Bystander tracks           & 40 \\
\textbf{Total tracks}      & \textbf{66} \\
\bottomrule
\end{tabular}
\end{singlespace}
\end{table}

\subsection{Classification Performance: Phase 2a (Geometry-Only)}
\label{sec:results_2a}

The Phase~2a configuration uses only the two bounding-box geometry cues
(area ratio and center distance) together with temporal hysteresis and latching.
No facial landmark extraction is performed. I evaluated the system against the
66 ground-truth tracks.

Table~\ref{tab:confusion_2a} presents the confusion matrix.

\begin{table}[htbp]
\begin{singlespace}
\centering
\caption{Confusion matrix for the Phase~2a (geometry-only) configuration. Rows give
the actual (human-labelled) class of each face track and columns give the system's
decision.}
\label{tab:confusion_2a}
\renewcommand{\arraystretch}{1.3}
\begin{tabular}{@{}lccc@{}}
\toprule
 & \textbf{Predicted Partner} & \textbf{Predicted Bystander} & \textbf{Total} \\
\midrule
\textbf{Actual partner}   & 22 (TP) & 4 (FN) & 26 \\
\textbf{Actual bystander} & 5 (FP)  & 35 (TN) & 40 \\
\textbf{Total}        & 27      & 39      & 66 \\
\bottomrule
\end{tabular}
\end{singlespace}
\end{table}

The five false positives represent bystanders that the geometry rule incorrectly
classified as partners (privacy leaks), and the four false negatives represent
interaction partners that were incorrectly blurred (utility losses).
Table~\ref{tab:metrics_2a} summarises all derived metrics.

\begin{table}[htbp]
\begin{singlespace}
\centering
\caption{Classification metrics for the Phase~2a (geometry-only) configuration.}
\label{tab:metrics_2a}
\renewcommand{\arraystretch}{1.3}
\begin{tabular}{@{}lcc@{}}
\toprule
\textbf{Metric} & \textbf{Value} & \textbf{Definition} \\
\midrule
Privacy Success Rate (PSR)          & 87.50\% & TN / (TN + FP) \\
Interaction Preservation Rate (IPR) & 84.62\% & TP / (TP + FN) \\
Precision                           & 81.48\% & TP / (TP + FP) \\
Recall                              & 84.62\% & TP / (TP + FN) \\
F$_1$ Score                         & 83.02\% & Harmonic mean of Precision and Recall \\
\bottomrule
\end{tabular}
\end{singlespace}
\end{table}

The Phase~2a system correctly blurred 35 of 40 bystander tracks (PSR = 87.50\%) and
correctly preserved 22 of 26 partner tracks (IPR = 84.62\%). The F$_1$ score of
83.02\% reflects a reasonably balanced trade-off between privacy protection and
interaction preservation.

\subsection{Classification Performance: Phase 2b (Full System with Landmark Cues)}
\label{sec:results_2b}

The Phase~2b configuration extends Phase~2a by adding the two facial landmark
cues derived from MediaPipe Face Landmarker: head-pose yaw (orientation veto) and
Mouth Aspect Ratio (V-VAD boost). I ran the complete Phase~2b pipeline on all
evaluation clips with the threshold values listed in Table~\ref{tab:params}.

\paragraph{Evaluation protocol.}
The 66-track ground truth of Section~4.2 could not be reused directly for Phase~2b.
SORT assigns track identifiers sequentially within each clip, so a single additional
or missing detection early in a clip shifts every subsequent identifier. This
occurred between the Phase~2a and Phase~2b processing runs, and the stored per-track
labels no longer referred to the same individuals. I therefore evaluated Phase~2b
through a blind, multi-annotator manual annotation study.

Ten evaluation frames were drawn automatically, one from each of the first five
bystander clips and the first five interaction clips. Within each clip, the frame
with the most clearly visible faces was selected (bounding box of at least
50~pixels per side and detection confidence of at least 0.50, excluding the first
30 frames while the tracker initialises). Each face in the selected frames was
marked with a plain numbered badge, with no colour or label indicating the system's
decision. Four annotators (myself and three independent annotators)
labelled every face as \emph{interaction partner} or \emph{bystander}, or left it
blank if undecidable; 114 of the 116 possible labels were provided. The final label
of each face was obtained by majority vote, with ties resolved as \emph{bystander}
(the privacy-conservative choice); one tie occurred. Each annotated face corresponds
to a SORT track, and, consistent with the latching protocol of Section~4.1, the
system's prediction for a face is \emph{partner} if its track was classified as
\textsc{Interactor} at any point in the clip.

I measured inter-annotator agreement with Cohen's~$\kappa$ for each of the six
annotator pairs. Pairwise $\kappa$ ranged from 0.307 to 0.724 (mean 0.587), with
five of the six pairs showing moderate or substantial agreement. The majority vote
yielded 29 annotated faces: 7 interaction partners and 22 bystanders
(Table~\ref{tab:sample_2b}).

\begin{table}[htbp]
\begin{singlespace}
\centering
\caption{Composition of the human-annotated Phase~2b evaluation sample.}
\label{tab:sample_2b}
\renewcommand{\arraystretch}{1.3}
\begin{tabular}{@{}lcccc@{}}
\toprule
\textbf{Clip category} & \textbf{Frames} & \textbf{Faces} & \textbf{Partners} & \textbf{Bystanders} \\
\midrule
Interaction & 5  & 15 & 7 & 8  \\
Bystander   & 5  & 14 & 0 & 14 \\
\midrule
\textbf{Total} & \textbf{10} & \textbf{29} & \textbf{7} & \textbf{22} \\
\bottomrule
\end{tabular}
\end{singlespace}
\end{table}

\paragraph{Results.}
Table~\ref{tab:confusion_2b} presents the confusion matrix for Phase~2b against the
majority-vote labels.

\begin{table}[htbp]
\begin{singlespace}
\centering
\caption{Confusion matrix for the Phase~2b (full system) configuration on the
human-annotated sample. Rows give the actual (human-labelled) class of each face and
columns give the system's decision.}
\label{tab:confusion_2b}
\renewcommand{\arraystretch}{1.3}
\begin{tabular}{@{}lccc@{}}
\toprule
 & \textbf{Predicted Partner} & \textbf{Predicted Bystander} & \textbf{Total} \\
\midrule
\textbf{Actual partner}   & 4 (TP) & 3 (FN)  & 7 \\
\textbf{Actual bystander} & 0 (FP) & 22 (TN) & 22 \\
\textbf{Total}        & 4      & 25      & 29 \\
\bottomrule
\end{tabular}
\end{singlespace}
\end{table}

\begin{table}[htbp]
\begin{singlespace}
\centering
\caption{Classification metrics for the Phase~2a and Phase~2b configurations.}
\label{tab:metrics_2b}
\renewcommand{\arraystretch}{1.3}
\begin{tabular}{@{}lcc@{}}
\toprule
\textbf{Metric} & \textbf{Phase 2a} & \textbf{Phase 2b (Full System)} \\
\midrule
Evaluation set                      & 66 tracks, 1 annotator & 29 faces, 4 annotators \\
\midrule
Privacy Success Rate (PSR)          & 87.50\% & 100.00\% \\
Interaction Preservation Rate (IPR) & 84.62\% & 57.14\% \\
Precision                           & 81.48\% & 100.00\% \\
Recall                              & 84.62\% & 57.14\% \\
F$_1$ Score                         & 83.02\% & 72.73\% \\
\bottomrule
\end{tabular}
\end{singlespace}
\end{table}

Phase~2b produced no false positives: all 22 bystanders in the sample were blurred
(PSR = 100\%), including the 8 bystanders who appeared in interaction clips
alongside an active partner. Of the 7 interaction partners, 4 were correctly
preserved and 3 were blurred (IPR = 57.14\%). Broken down by clip category, all 14
faces in bystander clips were correctly blurred, and in interaction clips the
system preserved 4 of 7 partners while blurring all 8 bystanders.

Because the two configurations were evaluated on different samples, the
comparison in Table~\ref{tab:metrics_2b} is indicative rather than conclusive. The
Phase~2b sample is small: with only seven partners, each misclassified partner
changes IPR by 14.3 percentage points. The 95\% Wilson confidence intervals are
85.1-100\% for Phase~2b PSR and 25.0-84.2\% for Phase~2b IPR, compared with
73.9-94.5\% (PSR) and 66.5-93.9\% (IPR) for Phase~2a; the intervals overlap for
both metrics. The causes of the three Phase~2b false negatives are analysed in
Section~\ref{sec:disc_failures}.

\subsection{Phase 2b: Video-Clip Evaluation}
\label{sec:results_clips}

The still-frame study had two weaknesses. Its sample was small (seven partners),
and annotators had to judge interaction from a single image, which explains its
moderate agreement. I therefore ran a second, larger blind evaluation of the same
Phase~2b system using short video clips with sound.

\paragraph{Evaluation protocol.}
I sampled face tracks from all 21 evaluation clips without reference to the
system's decisions: in each interaction clip, the four longest tracks plus two
randomly chosen tracks, and in each bystander clip, the four longest tracks. A track
was eligible if it was visible for at least 60 frames (about two seconds) and its
median face size was at least 32~pixels. This gave 104 clips (83 from interaction
videos and 21 from bystander videos). Each clip lasts seven seconds: a one-second
still frame showing which face to watch, followed by six seconds of the original,
unblurred video with its original sound. A single yellow box marked the target face,
and the clips were shown in a shuffled order with no indication of their source
category or of the system's decision.

For this study I defined an \emph{interaction partner} as any member of the camera
wearer's group or conversation. This includes a person talking to the wearer and
also a person who is with the wearer but not speaking at that moment, such as a
friend at the same table or playing the same game. A \emph{bystander} is anyone not
in the wearer's group, such as a stranger walking past or a person in a different
group. This is broader than the still-frame instruction (actively talking to the
wearer), and reflects the social reality that the members of the wearer's own group
are not the people whose privacy is at risk. Three annotators (myself and two
independent annotators) labelled each clip through a simple web page. The final
label was obtained by majority vote. Two clips were marked undecidable by two of the
three annotators and were excluded, leaving 102 labelled clips: 67 partners and 35
bystanders.

The system's decision for each clip is the decision visible in its output video: a
face counts as preserved if it was left unblurred in at least half of the frames in
which it was detected during the clip. I verified these decisions directly against
the rendered output videos, and all 104 agreed with the decisions recorded by the
pipeline.

\paragraph{Agreement.}
Agreement between annotators was almost perfect (Table~\ref{tab:clip_kappa}). The
pairwise Cohen's~$\kappa$ ranged from 0.913 to 0.957 and Fleiss'~$\kappa$ across all
three annotators was 0.942, compared with a mean pairwise $\kappa$ of 0.587 in the
still-frame study. Seeing movement and hearing speech therefore made the labelling
task far more reliable.

\begin{table}[htbp]
\begin{singlespace}
\centering
\caption{Inter-annotator agreement in the video-clip evaluation. Annotator A is
myself.}
\label{tab:clip_kappa}
\renewcommand{\arraystretch}{1.3}
\begin{tabular}{@{}lcc@{}}
\toprule
\textbf{Annotator pair} & \textbf{Cohen's $\kappa$} & \textbf{Raw agreement} \\
\midrule
A and B & 0.957 & 98.0\% \\
A and C & 0.913 & 96.1\% \\
B and C & 0.956 & 98.0\% \\
\midrule
\textbf{All three (Fleiss' $\kappa$)} & \textbf{0.942} & \\
\bottomrule
\end{tabular}
\end{singlespace}
\end{table}

\paragraph{Results.}
Table~\ref{tab:confusion_clips} gives the confusion matrix and
Table~\ref{tab:metrics_clips} the metrics, overall and by clip category.

\begin{table}[htbp]
\begin{singlespace}
\centering
\caption{Confusion matrix for Phase~2b in the video-clip evaluation. Rows give the
actual (human-labelled) class of each face and columns give the system's decision.}
\label{tab:confusion_clips}
\renewcommand{\arraystretch}{1.3}
\begin{tabular}{@{}lccc@{}}
\toprule
 & \textbf{Predicted Partner} & \textbf{Predicted Bystander} & \textbf{Total} \\
\midrule
\textbf{Actual partner}   & 38 (TP) & 29 (FN) & 67 \\
\textbf{Actual bystander} & 2 (FP)  & 33 (TN) & 35 \\
\textbf{Total}            & 40      & 62      & 102 \\
\bottomrule
\end{tabular}
\end{singlespace}
\end{table}

\begin{table}[htbp]
\begin{singlespace}
\centering
\caption{Phase~2b metrics in the video-clip evaluation, with 95\% Wilson confidence
intervals. IPR is not defined for bystander clips, which contain no partners.}
\label{tab:metrics_clips}
\renewcommand{\arraystretch}{1.3}
\begin{tabular}{@{}lccc@{}}
\toprule
\textbf{Metric} & \textbf{All clips} & \textbf{Interaction clips} & \textbf{Bystander clips} \\
\midrule
Faces evaluated & 102 & 81 & 21 \\
PSR       & 94.29\% (81.4-98.4) & 92.86\% (68.5-98.7) & 95.24\% (77.3-99.2) \\
IPR       & 56.72\% (44.8-67.9) & 56.72\% (44.8-67.9) & not defined \\
Precision & 95.00\% & 97.44\% & not defined \\
F$_1$     & 71.03\% & 71.70\% & not defined \\
\bottomrule
\end{tabular}
\end{singlespace}
\end{table}

The system blurred 33 of 35 bystanders (PSR = 94.29\%), including 20 of the 21 faces
in bystander clips and 13 of the 14 bystanders who appeared in interaction clips.
When the system left a face unblurred, it was correct 95\% of the time. It preserved
38 of 67 partners (IPR = 56.72\%). These figures closely match the still-frame study
(PSR 100\%, IPR 57.14\%), but the IPR confidence interval is now about half as wide
(44.8-67.9\% against 25.0-84.2\%), so the video-clip study confirms the earlier
finding with much stronger evidence. Of the 29 partners that were blurred, 24 had a
median face size below the 1.2\% area threshold, and 10 were confirmed as partners
by the system later in the same video, after the six-second clip had ended. These
failures are analysed in Section~\ref{sec:disc_failures}.

\subsection{Offline Second Pass for Recorded Video}
\label{sec:results_secondpass}

The interaction latch already ensures that a confirmed partner is never blurred
again. However, it works only forward in time: a partner is blurred until the
system has gathered enough evidence to confirm them. When the video is recorded and
processed after capture, the whole of each track is known in advance, so this
limitation can be removed. I implemented a second processing pass over the recorded
decisions, which applies two rules.

\begin{enumerate}
  \item \textbf{Retroactive latch.} A track that is confirmed as a partner at any
        point is preserved for its whole time on screen, including the frames before
        confirmation. This applies the principle that once a person is an interaction
        partner, they remain one.
  \item \textbf{Friend group.} A track that stays in view for at least ten seconds,
        and is on screen together with a confirmed partner for at least half of that
        time, is treated as a member of the wearer's group and preserved as well.
        This covers friends who sit with the wearer but do not speak, while strangers
        who pass by are rarely in view for that long. Requiring the presence of a
        confirmed partner prevents the rule from preserving strangers who simply stay
        in one place, such as a shop assistant.
\end{enumerate}

Neither rule requires additional model inference, since both operate on the
decisions the pipeline has already made. Table~\ref{tab:secondpass} reports their
effect on the 102 clips of the video-clip evaluation.

\begin{table}[htbp]
\begin{singlespace}
\centering
\caption{Effect of the offline second pass on the video-clip evaluation (102 clips,
95\% Wilson confidence intervals).}
\label{tab:secondpass}
\renewcommand{\arraystretch}{1.3}
\begin{tabular}{@{}lccc@{}}
\toprule
\textbf{Metric} & \textbf{Live system} & \textbf{+ Retroactive latch} & \textbf{+ Friend group} \\
\midrule
TP / FN / FP / TN & 38 / 29 / 2 / 33 & 48 / 19 / 3 / 32 & 53 / 14 / 3 / 32 \\
PSR       & 94.29\% (81.4-98.4) & 91.43\% (77.6-97.0) & 91.43\% (77.6-97.0) \\
IPR       & 56.72\% (44.8-67.9) & 71.64\% (59.9-81.0) & 79.10\% (67.9-87.1) \\
Precision & 95.00\% & 94.12\% & 94.64\% \\
F$_1$     & 71.03\% & 81.36\% & 86.18\% \\
\bottomrule
\end{tabular}
\end{singlespace}
\end{table}

The retroactive latch recovered 10 partners at the cost of one additional
bystander being preserved. The friend-group rule recovered a further 5 partners
without any additional privacy leak. Together, the two rules raised IPR from 56.72\%
to 79.10\% and F$_1$ from 71.03\% to 86.18\%, while PSR fell from 94.29\% to 91.43\%.
To check that the ten-second limit was not simply fitted to these clips, I repeated
the evaluation with limits from 10 to 60~seconds. Every limit in this range gave the
same PSR (91.43\%) and an IPR between 73.13\% and 79.10\%, so the result is not
sensitive to the exact value. Only a much shorter limit of five seconds admitted an
additional bystander.

Both rules were designed after examining the errors of the video-clip evaluation and
are evaluated on the same labels, so these figures should be read as an indication
of the benefit of the second pass rather than as an independent test. Both rules
also require the whole recording and therefore apply to recorded video, not to
live streaming.

\subsection{Comparison with Blanket-Blurring Baseline}
\label{sec:baseline_comparison}

The blanket-blurring baseline applies Gaussian blur to every detected face in every
frame, making no use of interaction cues. By construction, this baseline achieves a
PSR of 100\% (every bystander is blurred) and an IPR of 0\% (every interaction
partner is also blurred). The proposed system represents a qualitative departure
from the baseline: Phase~2a sacrifices a small amount of PSR to recover interaction
utility, while Phase~2b stays close to the baseline's PSR and still recovers a
large part of the interaction utility (Table~\ref{tab:vs_baseline}).

\begin{table}[htbp]
\begin{singlespace}
\centering
\caption{Comparison of the proposed system against the blanket-blurring baseline.
Phase~2a was evaluated on the 66-track ground truth, Phase~2b on the 29-face
still-frame sample (Section~\ref{sec:results_2b}) and on the 102-clip video sample
(Sections~\ref{sec:results_clips} and~\ref{sec:results_secondpass}).}
\label{tab:vs_baseline}
\renewcommand{\arraystretch}{1.3}
\begin{tabular}{@{}lccc@{}}
\toprule
\textbf{Configuration} & \textbf{PSR} & \textbf{IPR} & \textbf{F$_1$} \\
\midrule
Blanket-blur baseline      & 100\%   & 0\%     & N/A \\
Phase~2a (geometry)        & 87.50\% & 84.62\% & 83.02\% \\
Phase~2b, still frames     & 100.00\% & 57.14\% & 72.73\% \\
Phase~2b, video clips      & 94.29\% & 56.72\% & 71.03\% \\
Phase~2b + offline second pass & 91.43\% & 79.10\% & 86.18\% \\
\bottomrule
\end{tabular}
\end{singlespace}
\end{table}

The baseline's perfect PSR is achieved at the cost of complete utility destruction:
the camera wearer's interaction partners are indistinguishable from strangers in the
output stream. The proposed Phase~2a system recovers 84.62\% of interaction partners
(IPR) while correctly anonymizing 87.50\% of bystanders (PSR), at no additional
hardware cost relative to the baseline, since both configurations share the same
YOLOv8n-face detector and SORT tracker. In the video-clip evaluation, the Phase~2b
system blurred 94.29\% of bystanders while preserving 38 of 67 partners that the
baseline would have blurred, and with the offline second pass it preserved 53 of
67 partners (79.10\%) while still blurring 91.43\% of bystanders.

\subsection{Computational Performance}
\label{sec:fps}

I measured pipeline latency by re-running the full processing loop on a laptop
CPU (Intel Core i-series, no GPU or NPU acceleration enabled) and timing each stage
independently with a high-resolution wall-clock timer over a 300-frame clip. Table~\ref{tab:fps}
reports the per-stage mean latency for the Phase~2a configuration.

\begin{table}[htbp]
\begin{singlespace}
\centering
\caption{Per-stage mean latency (Phase~2a, 300-frame clip, laptop CPU).}
\label{tab:fps}
\renewcommand{\arraystretch}{1.3}
\begin{tabular}{@{}lcc@{}}
\toprule
\textbf{Stage} & \textbf{Mean latency (ms)} & \textbf{Fraction of total} \\
\midrule
Frame read          & 10.0  & 11.6\% \\
Face detection (YOLOv8n-face) & 55.1 & 64.1\% \\
Tracking (SORT)     & 20.1  & 23.4\% \\
Classification      & 0.02  & $<$0.1\% \\
Selective anonymization & 0.87  & 1.0\% \\
\midrule
\textbf{Total per frame} & \textbf{86.0} & \textbf{100\%} \\
\textbf{Throughput}      & \multicolumn{2}{c}{\textbf{11.62 FPS}} \\
\bottomrule
\end{tabular}
\end{singlespace}
\end{table}

Face detection with YOLOv8n-face dominates the pipeline at 55.1~ms per frame
(64.1\% of total latency). The classification stage is negligible at 0.02~ms,
confirming that the heuristic decision layer imposes no measurable overhead relative
to the underlying detection and tracking stages.

\paragraph{Raspberry Pi 5 projection.}
I did not benchmark directly on a Raspberry Pi 5, as the target hardware was not
available to me at the time of evaluation. Following the methodology described in
Section~3.8, a conservative hardware scaling factor of 3.5$\times$ is applied to
project performance on the ARM Cortex-A76 processor of the Raspberry Pi 5. This
factor is derived from published YOLO inference benchmarks comparing Intel Core-class
and ARM Cortex-A76 class processors on equivalent workloads. Applying this factor to
the measured 86.0~ms per-frame latency yields an estimated throughput of
approximately \textbf{3.3 FPS} on the Raspberry Pi 5 for the Phase~2a configuration.

\paragraph{Phase 2b landmark overhead.}
I benchmarked the landmark stage separately by timing every MediaPipe Face
Landmarker call within the same 300-frame loop, on the same laptop CPU under mains
power. Across two runs, landmark extraction measured \textbf{20.1~ms per face crop}
on average (21.1 and 19.1~ms); the benchmark clip averaged 0.99 tracked faces per
frame, adding 19.9~ms per frame. This is consistent with the 15-25~ms per-face
estimate anticipated in the methodology. Because detection latency varied between
runs owing to background CPU load, the Phase~2b figure is obtained by adding the
measured landmark cost to the Phase~2a baseline: $86.0 + 19.9 \approx 106$~ms per
frame, or approximately \textbf{9.4 FPS} on the laptop CPU (the fastest directly
measured Phase~2b run reached 9.9~FPS, in agreement). Applying the same
3.5$\times$ scaling factor gives approximately \textbf{2.7 FPS} on the Raspberry
Pi 5. Since landmark extraction runs once per tracked face, this overhead grows
linearly with the number of faces in view.

\paragraph{Battery power.}
Wearable devices always run on battery, and a laptop on battery limits its processor
speed to save energy. I therefore also benchmarked the pipeline with the laptop
unplugged. On battery, landmark extraction took
33.5~ms per face, and the complete Phase~2b pipeline measured 163.8~ms per frame
(6.1~FPS); two Phase~2a runs measured 9.1 and 7.5~FPS. These runs were also
affected by rising processor temperature over consecutive runs, so they represent a
conservative, worst-case condition rather than a precise figure.
Table~\ref{tab:power} summarises both conditions.

\begin{table}[htbp]
\begin{singlespace}
\centering
\caption{Measured laptop CPU throughput under mains and battery power, with the
projected Raspberry Pi 5 throughput (3.5$\times$ scaling factor). Battery figures are
single runs affected by rising processor temperature.}
\label{tab:power}
\renewcommand{\arraystretch}{1.3}
\begin{tabular}{@{}lcccc@{}}
\toprule
 & \multicolumn{2}{c}{\textbf{Mains power}} & \multicolumn{2}{c}{\textbf{Battery power}} \\
\textbf{Configuration} & \textbf{Laptop} & \textbf{Pi 5 (projected)} & \textbf{Laptop} & \textbf{Pi 5 (projected)} \\
\midrule
Phase~2a & 11.6~FPS & 3.3~FPS & 7.5-9.1~FPS & 2.1-2.6~FPS \\
Phase~2b & 9.4~FPS  & 2.7~FPS & 6.1~FPS     & 1.7~FPS \\
Landmarks per face & 20.1~ms & & 33.5~ms & \\
\bottomrule
\end{tabular}
\end{singlespace}
\end{table}

In summary, the complete pipeline runs entirely on a CPU, without a GPU, at
11.62~FPS (Phase~2a) and approximately 9.4~FPS (Phase~2b) on a laptop under mains
power, and at 6.1-9.1~FPS on battery, with projected Raspberry Pi 5 throughput of
1.7-3.3~FPS across both conditions. The interaction classifier
itself adds less than 0.1\% to the processing time, so the selective decision comes
at almost no extra cost over blanket blurring. Section~5.3 discusses optimisations
that can raise throughput further on edge hardware.

% ==============================================================================
%  CHAPTER 5 — DISCUSSION
% ==============================================================================
\section{Discussion}
\label{sec:discussion}

This chapter interprets the experimental results in the context of the research
questions posed in Chapter~1, examines the failure modes of the proposed system,
compares performance with the most relevant prior work, and identifies the
limitations of the current evaluation.

\subsection{Classification Performance and Research Question}
\label{sec:disc_perf}

The central research question asked whether a context-aware, multi-cue heuristic
operating on geometry and facial landmarks could selectively anonymize bystanders
while preserving interaction partners, without biometric recognition or cloud
offloading. The Phase~2a evaluation results (PSR = 87.50\%, IPR = 84.62\%) confirm
that this is feasible using only bounding-box geometry. A PSR of 87.50\% means that
the system incorrectly exposes 12.5\% of bystander tracks; equivalently, 5 of the
40 bystander tracks in the evaluation dataset were mistakenly preserved unblurred.
An IPR of 84.62\% means that 15.38\% of interaction partners were unnecessarily
blurred; equivalently, 4 of 26 partner tracks were incorrectly anonymized.

Both error types carry concrete costs. A false positive (bystander preserved as
partner) constitutes a \emph{privacy leak}: the system fails to protect an
individual who did not consent to being recorded. A false negative (partner blurred
as bystander) constitutes a \emph{utility loss}: the wearer's conversation partner
is anonymized, defeating the purpose of selective anonymization. The observed ratio
of 5 FPs to 4 FNs reflects a roughly balanced trade-off, consistent with the
system's design objective of avoiding strong asymmetry in either direction.

The addition of Phase~2b landmark cues is expected to reduce false positives
specifically: the yaw veto correctly classifies faces that are turned significantly
away from the camera as bystanders, even if those faces happen to be large or
central in the frame (which would incorrectly pass the Phase~2a geometry rule). The
V-VAD boost conversely reduces false negatives by promoting faces that exhibit
sustained mouth movement to \emph{Interactor} status, even when centrality is
marginally insufficient.

The Phase~2b results (Table~\ref{tab:metrics_2b}) are consistent with the first
expectation: every bystander in the human-annotated sample was blurred, so no
bystander's face was left visible in the output (FP = 0, PSR = 100\%). This includes
eight bystanders who appeared in interaction clips next to an active partner. The
second expectation was not borne out: IPR was 57.14\%, with three of seven partners
blurred. As Section~\ref{sec:disc_failures} shows, these partners were not blurred
because the head-pose check judged them to be facing away. They were blurred because
their faces were too small in the frame to pass the proximity check, which is applied
before the head-pose and speaking checks are considered. In the
evaluated sample, the full system therefore behaved as a privacy-first classifier.

The larger video-clip evaluation confirms this pattern with much stronger evidence.
With almost perfect agreement between annotators ($\kappa = 0.942$) and 102 clips,
the live Phase~2b system blurred 94.29\% of bystanders and preserved 56.72\% of
partners, and when it preserved a face it was correct 95\% of the time. The
privacy-first behaviour of the live system is therefore a stable property rather
than an artefact of a small sample. The offline second pass shows that most of the
remaining utility loss can be recovered without weakening privacy substantially:
by applying the latch retroactively and recognising the wearer's group, IPR rose to
79.10\% at a PSR of 91.43\%, giving the best balance of the configurations I
evaluated (F$_1$ = 86.18\%). Together, these results answer the research question
positively: interaction-aware selective anonymization can protect most bystanders
while preserving most of the people the wearer is with, without face recognition or
additional sensors.

\subsection{Comparison with BystandAR}
\label{sec:disc_bystandar}

The closest prior system, BystandAR \cite{bystandar}, achieved PSR = 98.14\% and
IPR = 96.27\% on a 16-participant evaluation with dedicated gaze-tracking and
directional microphone hardware. The proposed system achieves lower accuracy (PSR
87.50\%, IPR 84.62\% for Phase~2a; PSR 100\%, IPR 57.14\% for Phase~2b on a smaller
annotated sample) using only a monocular RGB camera and CPU-only inference. This
performance gap is expected and attributable to two factors: the richness of input
modalities and the evaluation scale.

BystandAR uses eye gaze as a direct signal of where the wearer is looking, which is
a near-perfect proxy for interaction intent. The proposed system infers this
indirectly from face geometry and pose, which is a weaker signal subject to ambiguity
when bystanders happen to be close to the camera wearer (for example, in a crowded
corridor). Additionally, BystandAR's evaluation involved 16 controlled participant
sessions with standardised interaction protocols, whereas the present evaluation used
21 naturalistic clips without controlled conditions.

Despite this gap, the proposed system's significance lies in its deployment
constraints: it requires no gaze tracker, no directional microphone array, and no
dedicated wearable hardware. It operates on a commodity RGB camera using only
CPU-bound inference, making it directly deployable on low-cost edge devices without
hardware modification. This positions it as a practical privacy layer for consumer
AR glasses such as the Meta Ray-Ban, which lack the dedicated sensors required by
BystandAR.

\subsection{Computational Feasibility and Optimisation Pathways}
\label{sec:disc_fps}

The measured throughput of 11.62 FPS on a laptop CPU shows that the complete
pipeline runs at interactive rates without any GPU acceleration. On the Raspberry Pi 5,
the projected throughput is approximately 3.3 FPS (Phase~2a) and 2.7 FPS (Phase~2b)
with the unoptimised PyTorch implementation. Since face detection accounts for 64\% of
the processing time, three optimisation pathways can raise this substantially:

\begin{enumerate}
  \item \textbf{ONNX export and quantisation.} Exporting YOLOv8n-face to ONNX INT8
        format and running inference via ONNXRuntime with NEON-optimised kernels for
        ARM can reduce detection latency by a factor of 2-3$\times$ on ARM hardware
        relative to the PyTorch runtime used in the current evaluation.

  \item \textbf{Reduced input resolution.} Processing frames at $416 \times 416$
        rather than $640 \times 640$ reduces YOLO inference time by approximately
        35\% with a modest reduction in detection recall for small, distant faces.
        For the target use case of close-range egocentric video, this trade-off is
        likely acceptable.

  \item \textbf{Landmark extraction batching.} In the current implementation,
        MediaPipe Face Landmarker is called once per face per frame. Batching
        multiple face crops into a single inference call would reduce per-face
        overhead, particularly when multiple faces are present simultaneously.
\end{enumerate}

With all three optimisations applied, achieving 8-12~FPS on the Raspberry Pi 5 is
plausible, which is sufficient for processing recorded egocentric video (a common
use case in AR privacy logging). Full-frame-rate operation on the Pi~5 can be reached by
running the detector on a neural processing unit (NPU) or AI co-processor, which
is included in many current smart-glasses processors.

\paragraph{Comparison with commercial smart-glasses hardware.}
To judge whether the evaluation hardware is representative, I compared it with
current camera-equipped glasses (Table~\ref{tab:glasses}). Ray-Ban Meta and Meta
Ray-Ban Display glasses use the Qualcomm Snapdragon AR1 Gen~1 platform
\cite{ref:qcom_ar1}, which combines a quad-core CPU (reported at up to about
1.9~GHz), dual image signal processors and a dedicated Hexagon neural processing unit
(NPU), with 2~GB of RAM in the Display model \cite{ref:gsmarena_display}. Snap
Spectacles (fifth generation) split the work across two Qualcomm Snapdragon processors
and run for about 45 minutes per charge \cite{ref:spectacles5}, and Meta's Aria
Gen~2 research glasses perform their machine perception on a custom low-power
co-processor \cite{ref:aria2}.

\begin{table}[htbp]
\begin{singlespace}
\centering
\caption{Evaluation hardware compared with current camera-equipped glasses.
Values are from manufacturer and press sources; ``not published'' means the
manufacturer has not released the figure.}
\label{tab:glasses}
\renewcommand{\arraystretch}{1.3}
\small
\begin{tabular}{@{}lllll@{}}
\toprule
\textbf{Device} & \textbf{Processor} & \textbf{AI accelerator} & \textbf{Battery} & \textbf{Where vision AI runs} \\
\midrule
Laptop (my evaluation)    & Intel Core i5-13420H & none used & mains / battery & on device (CPU) \\
Raspberry Pi 5 (target) & 4$\times$ Cortex-A76, 2.4~GHz & none & external & on device (CPU) \\
Ray-Ban Meta          & Snapdragon AR1 Gen~1 & Hexagon NPU & about 4~h & mostly phone and cloud \\
Meta Ray-Ban Display  & Snapdragon AR1 Gen~1 & Hexagon NPU & about 6~h & mostly phone and cloud \\
Snap Spectacles (5th gen.) & 2$\times$ Snapdragon & not published & about 45~min & on device \\
Meta Aria Gen~2       & custom co-processor & custom & 6-8~h & on device \\
\bottomrule
\end{tabular}
\end{singlespace}
\end{table}

This comparison supports the design choices of my research in three ways. First, the
Raspberry Pi 5 is a reasonable stand-in for the CPU of current glasses: its four
Cortex-A76 cores at 2.4~GHz are of the same class as the AR1's quad-core CPU, so
the projected Pi~5 figures are a realistic estimate of CPU-only performance on such
devices. Second, my evaluation used the CPU only, whereas every current glasses
platform includes a neural processing unit. Face detection, which accounts for 64\%
of the pipeline's processing time, is exactly the kind of workload an NPU
accelerates. On a Snapdragon smartphone NPU, a quantised YOLO detector runs in about
5.6~ms per frame against 53~ms on the same phone's CPU \cite{ref:ultralytics_qnn},
and Qualcomm's quantised face detector runs in under 1~ms \cite{ref:qcom_facedet}.
Glasses processors are smaller than smartphone processors, so these figures are not
directly transferable, but they indicate that the CPU-only results reported here are
a worst case for real glasses. Third, battery capacity, rather than raw speed, is
the main constraint on glasses: Spectacles run for about 45 minutes, and current
Ray-Ban Meta glasses send most vision AI to the paired phone or the cloud. Because
my pipeline keeps all face processing on the device, it avoids sending bystanders'
faces off the glasses, which Meta itself identifies as a central privacy concern for
camera glasses \cite{ref:meta_bystander}. Reducing how often the detector runs, as
proposed above, also directly reduces energy use.

\subsection{Failure Mode Analysis}
\label{sec:disc_failures}

Analysis of the five false-positive tracks (bystanders incorrectly preserved as
partners) reveals a consistent pattern: all five occurred in the \emph{interaction}
video category, where people standing near the camera wearer were never directly
addressed but occupied a large, central position in the frame for an extended
period. In these cases, the area ratio and centrality thresholds are satisfied
because the individual is physically close to the wearer, even though no interaction
occurs. This ambiguity is precisely the scenario that the Phase~2b yaw veto is
designed to address: a person who is nearby but looking in a different direction
(yaw $> 40^\circ$) should be vetoed as a bystander regardless of proximity.

The four false-negative tracks (partners incorrectly blurred) arose from brief
interactions that began after a long period in which the face was peripheral, so
the temporal hysteresis had not yet committed the track to the \emph{Interactor}
state.

\paragraph{Phase 2b failure modes.}
All three Phase~2b false negatives share a single cause: at the annotated frame,
each partner's face occupied between 0.36\% and 0.72\% of the frame area, below the
area threshold $\tau_{\mathrm{area}} = 1.2\%$. Two of them occur in one frame of a
seated group setting, where the same person was covered by two overlapping
detections that SORT tracked as separate tracks; this duplicate-track artefact
counts one missed partner twice. The third is a partner standing at a service
counter near the edge of the frame, small and partly turned (yaw $\approx 25^\circ$,
below the $40^\circ$ veto). In each case the area gate rejected the face before any
landmark cue was considered. Because the V-VAD boost also requires
$r_{\mathrm{area}} \geq \tau_{\mathrm{area}}$, it cannot rescue partners who converse
from further away than the area threshold assumes, for example across a table.

A further factor limits the influence of the landmark cues. MediaPipe Face
Landmarker returned landmarks for only 50.8\% of the 224{,}519 logged face
detections. For the remainder, which are typically small, motion-blurred or partially
occluded faces, no yaw or MAR value is available and the classifier falls back to the
Phase~2a geometry rule. Both observations point to the same improvement: relaxing
the area threshold when landmarks confirm a frontal, speaking face, and improving
landmark coverage on small faces (for example by upscaling the crop before landmark
extraction).

\paragraph{Video-clip failure modes.}
The video-clip evaluation confirms the same dominant cause on a larger sample. Of
the 29 partners the live system blurred, 24 had a median face size below the area
threshold. The median face size of missed partners was 0.92\% of the frame, against
1.71\% for partners the system preserved. Most of these were friends sitting further
from the wearer around a table, who are partners under the group definition but
appear smaller than the threshold assumes. In addition, 10 of the 29 were confirmed
as partners later in the same video, which shows that the live system often
recognises a partner correctly but too late. The offline second pass addresses both
causes directly, and recovered 15 of the 29. Of the two bystanders the live system
preserved, one was left unblurred for exactly half of the clip, the threshold of the
per-clip decision rule, and the other was a person close to the wearer in an
interaction clip.

\subsection{Limitations}
\label{sec:disc_limitations}

The following limitations should be considered when interpreting the results of this
evaluation.

\textbf{Dataset scale and diversity.} The evaluation dataset comprises 21 clips
collected from online sources rather than recorded for this study, so capture
conditions such as camera model, mounting position and frame rate are uncontrolled
and vary between clips. Interaction was evaluated only indoors, and the interaction
clips are dominated by seated tabletop scenes; partner detection was therefore not
tested in outdoor conversations, while bystander behaviour was tested only outdoors.
Results may not generalise to other settings, lighting conditions or interaction
styles. A larger dataset in which both categories span the same range of
environments would provide a more robust assessment of generalisation.

\textbf{Annotation scale and protocol.} I assigned the Phase~2a ground truth
alone. Phase~2b was evaluated twice with blind protocols: on 29 faces from still
frames by four annotators (mean $\kappa = 0.587$), and on 102 video clips by three
annotators (Fleiss' $\kappa = 0.942$). I was one of the annotators in both studies
and was familiar with the system design, although my labels agreed with the
independent annotators' labels in 96-98\% of the video clips. The two Phase~2b
studies also used different partner definitions (actively talking to the wearer,
and member of the wearer's group). Because Phase~2a and Phase~2b were evaluated on
different samples, differences between them are indicative rather than conclusive.

\textbf{Offline second pass.} The retroactive latch and the friend-group rule were
designed after examining the errors of the video-clip evaluation and were evaluated
on the same labels. Although the result was stable across a wide range of time
limits, a fresh set of labelled clips is needed for an independent test. Both rules
also need the whole recording, so they apply to recorded video only.

\textbf{Heuristic thresholds.} I set all eight threshold parameters
(Table~\ref{tab:params}) empirically on the evaluation dataset itself. This risks
overfitting to the specific conditions of the collected clips. A held-out validation set or
cross-validation procedure would provide a fairer estimate of threshold robustness.

\textbf{Absence of audio.} The V-VAD cue relies entirely on lip movement as a
speech proxy and cannot distinguish between a person mouthing words, laughing,
chewing, or speaking. Incorporating even a rudimentary microphone-based voice
activity detector as a secondary cue would substantially reduce V-VAD ambiguity at
negligible hardware cost on most wearable devices.

\textbf{Hardware projection.} The Raspberry Pi 5 throughput figures (approximately
3.3 FPS for Phase~2a and 2.7 FPS for Phase~2b) are estimates derived from a scaling factor rather than a direct
measurement. They should be treated as an order-of-magnitude projection; actual
performance may differ depending on OS configuration, thermal throttling, and the
version of underlying inference libraries installed. The battery-power measurements
are single runs affected by rising processor temperature, and I did not measure
energy consumption directly, so the battery cost of continuous operation on glasses
remains to be quantified.

% ==============================================================================
%  CHAPTER 6 — CONCLUSION
% ==============================================================================
\section{Conclusion}
\label{sec:conclusion}

\subsection{Summary of Contributions}
\label{sec:conclusion_summary}

In this dissertation, I presented a context-aware, real-time privacy framework for
egocentric video that selectively anonymizes bystanders while preserving the faces
of active interaction partners. The framework operates as a five-stage sequential
pipeline: face detection via YOLOv8n-face, multi-face tracking via SORT, interaction
cue extraction, heuristic classification, and selective Gaussian blur anonymization.
All stages execute on CPU-only hardware without GPU acceleration, biometric face
recognition, or cloud offloading.

The primary contribution of my research is the \emph{multi-cue interaction
classifier}, which fuses four complementary signals derived from two information
sources. The first two cues (face area ratio and center distance) are extracted
from bounding-box geometry alone and require no additional model inference. The
second two cues (head-pose yaw via PnP and Mouth Aspect Ratio via MediaPipe Face
Landmarker) require facial landmark extraction but provide qualitatively different
information: yaw provides a hard veto for faces that are turned away from the
camera, while MAR-based V-VAD provides a temporal boost for faces exhibiting
sustained lip movement. The classifier combines these cues through a priority-ordered
decision function with temporal hysteresis and interaction latching, achieving PSR
= 87.50\% and IPR = 84.62\% in the Phase~2a geometry-only configuration, evaluated
against a ground-truth dataset of 66 unique face tracks from 21 egocentric video clips.
With landmark cues enabled (Phase~2b), a blind four-annotator evaluation of 29 faces
gave PSR = 100\% and IPR = 57.14\%, and a larger blind evaluation of 102 video
clips with almost perfect annotator agreement ($\kappa = 0.942$) confirmed this
with PSR = 94.29\% and IPR = 56.72\%. For recorded video, an offline second pass
that applies the latch retroactively and recognises the members of the wearer's
group raised IPR to 79.10\% at PSR = 91.43\% (F$_1$ = 86.18\%).

A secondary contribution is the \emph{systematic comparison with the blanket-blurring
baseline}. The baseline achieves PSR = 100\% but IPR = 0\%, making it unsuitable for
any use case in which the wearer needs to interact with people in the scene. The
proposed system recovers 84.62\% of interaction-partner utility (IPR) at a cost of
12.5 percentage points of PSR relative to the baseline. This trade-off is
appropriate for a system whose primary objective is to protect bystanders without
degrading the wearer's social experience. With the offline second pass, the
system preserved 79.10\% of the people in the wearer's group while still blurring
91.43\% of bystanders, demonstrating that selective anonymization can recover most
of the social value that blanket blurring destroys while keeping strong bystander
protection.

A third contribution is the \emph{edge-deployable implementation}. The pipeline is
designed from the outset to run on commodity CPU hardware. The measured throughput
of 11.62~FPS (Phase~2a) and approximately 9.4~FPS (Phase~2b) on a laptop CPU, with
estimated Raspberry Pi~5 throughput of 3.3 and 2.7~FPS respectively,
demonstrates that the approach is feasible on CPU-only edge hardware, and the
per-stage analysis identifies a clear path to higher frame rates through model
quantisation and hardware acceleration.

Finally, the \emph{evaluation methodology} that I developed for this project,
which includes blind annotation of still frames and of short video clips with
sound, inter-annotator agreement measurement and verification of the system's
decisions against its rendered output, forms a reusable protocol for future
egocentric privacy research.

\subsection{Relation to Prior Work}
\label{sec:conclusion_prior}

Compared to the most relevant prior system, BystandAR \cite{bystandar} (PSR~98.14\%,
IPR~96.27\%), the proposed framework achieves lower accuracy but eliminates the
hardware dependency on eye gaze trackers and directional microphones. Compared to
blanket-blurring systems such as EgoBlur \cite{egoblur2023} and SafeAR
\cite{ref:safear2025}, the proposed framework introduces the context-aware
classification layer that all prior blanket systems lack.

To my knowledge, the proposed system is the first to combine:
(i) real-time selective anonymization in egocentric video,
(ii) CPU-only edge deployment without biometric recognition,
(iii) a multi-cue fusion of geometry, head pose, and V-VAD, and
(iv) a quantitative evaluation distinguishing PSR from IPR as separate privacy and
utility metrics.

\subsection{Future Work}
\label{sec:conclusion_future}

I identify several directions for future investigation.

\textbf{Learned classifier.} The heuristic decision function could be replaced with
a lightweight learned binary classifier (e.g., a shallow MLP or logistic regression
model trained on the cue vectors) that generalises better across environments and
interaction styles than hand-tuned thresholds.

\textbf{Audio-visual fusion.} Incorporating a microphone-based voice activity
detector alongside the visual V-VAD cue would resolve ambiguous lip-movement cases
and substantially improve IPR in noisy environments.

\textbf{Independent and larger-scale evaluation.} The video-clip protocol should
be applied to a fresh set of clips, to test the offline second pass independently,
and to a larger, more diverse dataset, so that all configurations can be compared on
a common ground truth.

\textbf{Group recognition in live operation.} The retroactive latch and the
friend-group rule currently need the whole recording. Delaying the output stream
by a few seconds would give the system time to confirm partners and group members
before their frames are shown, bringing most of the benefit of the second pass to
live operation.

\textbf{Adaptive proximity threshold.} Most Phase~2b false negatives fell below the
fixed area threshold. Relaxing $\tau_{\mathrm{area}}$ when landmark cues confirm a
frontal, speaking face, and improving landmark coverage on small faces, are
direct routes to higher IPR without weakening PSR.

\textbf{On-device deployment.} Direct benchmarking on the Raspberry Pi~5 with ONNX
INT8 quantised models, combined with the input-resolution and batching optimisations
described in Section~5.3, would validate or refine the throughput projection and
identify any architecture-specific performance bottlenecks not captured by the
current scaling-factor estimate.

\textbf{Consent and transparency mechanisms.} The current system operates
invisibly from the bystander's perspective. Future work could integrate BystandARIA's
visible-light communication approach \cite{ref:bystandaria} as a complementary
consent channel, allowing bystanders to actively signal their privacy preferences
to the system.

% ==============================================================================
%  NEW REFERENCES: paste into \begin{thebibliography} in the main document
% ==============================================================================
% \bibitem{ref:qcom_ar1}
% Qualcomm Technologies, ``Snapdragon AR1 Gen 1 Platform,'' product page. [Online].
% Available: \url{https://www.qualcomm.com/products/mobile/snapdragon/xr-vr-ar/snapdragon-ar1-gen-1-platform}
%
% \bibitem{ref:gsmarena_display}
% GSMArena, ``Meta unveils Ray-Ban Display and Ray-Ban Meta (Gen 2) smart glasses,'' 2025.
% [Online]. Available: \url{https://m.gsmarena.com/meta_rayban_display_and_rayban_meta_gen_2_smart_glasses_debut-amp-69560.php}
%
% \bibitem{ref:spectacles5}
% Snap Inc., ``SPS 2024: Introducing the new Spectacles and Snap OS,'' Snap Newsroom,
% Sep. 2024. [Online]. Available: \url{https://newsroom.snap.com/sps-2024-spectacles-snapos}
%
% \bibitem{ref:aria2}
% Meta, ``Introducing Aria Gen 2: Unlocking new research in machine perception,
% contextual AI, robotics, and more,'' 2025. [Online]. Available:
% \url{https://www.meta.com/blog/project-aria-gen-2-next-generation-egocentric-research-glasses-reality-labs-ai-robotics/}
%
% \bibitem{ref:ultralytics_qnn}
% Ultralytics, ``Qualcomm QNN export for Ultralytics YOLO models,'' documentation.
% [Online]. Available: \url{https://docs.ultralytics.com/integrations/qnn}
%
% \bibitem{ref:qcom_facedet}
% Qualcomm AI Hub, ``Lightweight-Face-Detection-Quantized,'' model card, Hugging Face.
% [Online]. Available: \url{https://huggingface.co/qualcomm/Lightweight-Face-Detection-Quantized}
%
% \bibitem{ref:meta_bystander}
% Meta, ``Bystander privacy,'' white paper, 2026. [Online]. Available:
% \url{https://about.fb.com/wp-content/uploads/2026/07/Bystander-Privacy.pdf}
